\documentclass[11pt,a4paper]{article}
\usepackage[T1]{fontenc}
\usepackage{lmodern}
\usepackage{amsmath,amssymb}
\usepackage{booktabs}
\usepackage{array}
\usepackage[a4paper,margin=2cm]{geometry}
\usepackage[hidelinks]{hyperref}
\setlength{\parskip}{2pt}
\title{\vspace{-1.6cm}\Large Field-specific limits: phason relaxation and the magnetoelastic mixture}
\author{Dr.\ Vipin\\ Gurugram University}
\date{30 October 2026}
\begin{document}
\maketitle
\vspace{-6mm}

\section*{Neha: phason relaxation}

Bulk, decoupled: $K_1k^{2}=\rho_p\omega^{2}+\mathrm{i}\omega\gamma$ gives
$c_{ph}=\sqrt{K_1/\rho_p}=31622.78$~m/s, $k_i/k_r=\gamma/(2\rho_p\omega)$ below
$2\rho_p\omega=1.5\times10^{9}$, and the diffusive branch $k_r=\sqrt{\omega\gamma/(2K_1)}$,
$k_i/k_r\to1$ above it. At $\omega=2\pi10^{8}$:

\begin{center}\footnotesize
\begin{tabular}{lccc}
\toprule
$\gamma$ & 0 & $10^{8}$ & $10^{12}$\\
\midrule
$\omega/k_r$ (m/s) & 31622.8 & 31553.8 & 1228\\
$k_i/k_r$ & 0 & $6.60\times10^{-2}$ & 0.9993\\
closed form & 31622.8 & wave & $5.12\times10^{5}$\\
\bottomrule
\end{tabular}
\end{center}

Plate, all couplings on, $l=0$, $h=5\times10^{-4}$~m, $fd=1000$, A0: 2216.6~m/s at
$\gamma=10^{2},10^{6},10^{8}$, 2219.3 at $10^{10}$ and 2222.0 at $10^{12}$, against 2222.5 for the
decoupled phason. The phonon branch recovers the decoupled value once the phason is overdamped, and
the coupling costs 5.9~m/s while it still propagates. The Neha S0 branch sits near 5050~m/s with a
dip depth of 12, so I quote no S0 entry.

\section*{Anu: magnetoelastic coefficient against volume fraction}

The mixing rule is ours, not cited: $\alpha(v)=v(1-v)\alpha_0$, put in by scaling the three
magnetoelastic constants with $f=4v(1-v)$, so $\alpha^{2}\propto f^{2}$. Thin plate, $fd=100$: S0 is
6280.67~m/s for every $v$ to 0.01~m/s, because at small $k$ the magnetoelastic terms enter as
$\eta q_{31}$, $\eta q_{33}$, $\mathrm{i}kq_{15}$ and $q_{15}/c_{44}\sim4\times10^{-9}$.

A0 at $fd=1000$, window of $0.5$~m/s:

\begin{center}\footnotesize
\begin{tabular}{lccccc}
\toprule
$v$ & 0 & 0.25 & 0.5 & 0.75 & 1\\
\midrule
$f$ & 0 & 0.75 & 1 & 0.75 & 0\\
A0 (m/s) & 2276.17 & 2276.94 & 2277.48 & 2276.94 & 2276.17\\
$\Delta c$ (m/s) & 0 & $+0.775$ & $+1.310$ & $+0.775$ & 0\\
$\Delta c/f^{2}$ & --- & 1.377 & 1.310 & 1.377 & ---\\
\bottomrule
\end{tabular}
\end{center}

Peak at $v=1/2$, zero at both ends, $\Delta c/f^{2}$ constant to five per cent, which is the end of
the quadratic regime rather than scatter: the symmetric points agree to the last digit.

\section*{Parity and one note}

With all couplings on the families stay separate (Neha A0 2216.6, S0 near 5050; Anu A0 2277.1, S0
6280.7) and both papers keep the single pattern $(+,-,-,-,+)$ and its negative. For the scholar
file: $n_0^{th}=\epsilon_{33}\lvert-\mathrm{i}\omega+D_nk^{2}+1/\tau_R\rvert/(qb_n)$, so the
screening density follows the carrier lifetime and the 1~ns in the papers is a placeholder.

\end{document}
